\documentclass[11pt]{article}
\usepackage[utf8]{inputenc}
\usepackage[T1]{fontenc}
\usepackage{amsmath,amssymb,mathtools}
\usepackage{geometry}
\usepackage{booktabs}
\geometry{margin=1in}

\title{Depth-Shortcut Training Objective with Centered Pair Sampling and Logit-Normal Timesteps}
\author{}
\date{}

\begin{document}
\maketitle

\section{Active configuration}

This document describes the training path for the command using
\texttt{hybrid\_deep\_10}, direction--magnitude shortcut prediction,
centered pair sampling, logit-normal timestep sampling, output distillation,
and private activation regularization.

\begin{center}
\begin{tabular}{ll}
\toprule
Switch & Value \\
\midrule
\texttt{--shortcut-predictor} & \texttt{hybrid\_deep\_10} \\
\texttt{--shortcut-training-mode} & \texttt{direction-magnitude} \\
\texttt{--shortcut-predictor-use-timestep} & on \\
\texttt{--shortcut-predictor-use-class-input} & on \\
\texttt{--shortcut-predictor-class-fusion} & \texttt{add} \\
\texttt{--no-shortcut-predictor-normalize-input} & raw source activations \\
\texttt{--timestep-sampling-mode} & \texttt{logit\_normal} \\
\texttt{--direct-pair-mode} & \texttt{trunc\_normal\_centered} \\
\texttt{--output-distill-pair-mode} & \texttt{trunc\_normal\_centered} \\
\texttt{--pair-center-sigma} & $2.0$ \\
\texttt{--direct-num-pairs} & $1$ \\
\texttt{--output-distill-ratio} & $0.10$ \\
\texttt{--output-distill-update-mode} & \texttt{predictor\_plus\_all} \\
\texttt{--private-loss} & on \\
\texttt{--private-max-pairs} & $4$ \\
\texttt{--private-cosine-mode} & \texttt{bnd} \\
\texttt{--private-pair-mode} & \texttt{random} \\
\bottomrule
\end{tabular}
\end{center}

This run is a depth-shortcut run. It does not enable the timestep-shortcut
paired-batch objective and it does not enable depth-time resync.

\section{Timestep sampling}

Let the number of shortcut timestep bins be $T=50$. For each training example,
sample
\[
  \epsilon \sim \mathcal{N}(0,1),
  \qquad
  u=\sigma(\epsilon),
\]
where $\sigma$ is the logistic sigmoid. The discrete timestep index is
\[
  q = \operatorname{round}\bigl((T-1)u\bigr),
  \qquad
  q\in\{0,\ldots,T-1\},
\]
and
\[
  \tau = \frac{q}{T-1}.
\]
Thus the training timestep mass is concentrated near $\tau=0.5$ while still
covering both noisy and clean regions.

The repo convention is
\[
  \tau=0 \text{ is pure noise}, \qquad \tau=1 \text{ is clean data}.
\]
For a clean latent $x_0$ and Gaussian noise $x_1\sim\mathcal{N}(0,I)$,
\[
  x_\tau = (1-\tau)x_1 + \tau x_0.
\]
The flow-matching velocity target is
\[
  v^\star = x_0 - x_1.
\]

\section{Backbone forward}

The DiT backbone is evaluated once on the sampled noisy latent:
\[
  \left(
    \hat v_\theta,\{Z_\ell\}_{\ell=0}^{L},e_\tau
  \right)
  =
  F_\theta^{\mathrm{hidden}}(x_\tau,\tau,y),
\]
where $Z_0$ is the patch embedding and $Z_\ell$ for $\ell\ge1$ are block
outputs. For DiT-B, $L=12$.

The base velocity loss is
\[
  \mathcal{L}_{\mathrm{gen}}
  =
  \mathbb{E}\left[
    \left\|\hat v_\theta - v^\star\right\|_2^2
  \right].
\]

\section{Centered layer-pair sampling}

A candidate depth pair is
\[
  (a,b), \qquad 0\le a < b \le L.
\]
Define the layer gap and midpoint
\[
  d=b-a,\qquad m=\frac{a+b}{2}.
\]
The command uses
\[
  d_{\max}=10,\qquad \mu_d=3,\qquad \sigma_d=2,
  \qquad \mu_m=\frac{L}{2}=6,\qquad \sigma_m=2.
\]
Valid candidates satisfy $1\le d\le d_{\max}$.

The gap weight is a truncated discrete normal:
\[
  w_{\mathrm{gap}}(d)
  =
  \exp\left[
    -\frac{1}{2}
    \left(\frac{d-\mu_d}{\sigma_d}\right)^2
  \right].
\]
Within each gap, midpoint weights concentrate pairs around the middle of the
network:
\[
  w_{\mathrm{center}}(a,b)
  =
  \exp\left[
    -\frac{1}{2}
    \left(\frac{m-\mu_m}{\sigma_m}\right)^2
  \right].
\]
The implementation normalizes midpoint weights within each gap. Equivalently,
the pair distribution is proportional to
\[
  p(a,b)
  \propto
  w_{\mathrm{gap}}(b-a)
  \frac{
    w_{\mathrm{center}}(a,b)
  }{
    \sum_{a':\,a'+d\le L} w_{\mathrm{center}}(a',a'+d)
  }.
\]

\section{Shortcut predictor}

The predictor receives raw source activations because
\texttt{--no-shortcut-predictor-normalize-input} is set. For a sampled pair
$(a,b)$:
\[
  S_a = Z_a.
\]
It also receives the source log magnitude
\[
  m_a = \log\left(\|Z_a\|_2+\varepsilon\right),
\]
the current timestep embedding $e_\tau$, the layer ids $(a,b)$, and class label
conditioning. The predictor outputs a direction-like tensor and a normalized
log-magnitude delta:
\[
  (Y_{ab},\Delta m_{ab})
  =
  P_\phi(S_a,a,b,e_\tau,m_a,y).
\]

\section{Direct depth-shortcut loss}

For the sampled direct pair $(a,b)$, define the predicted direction
\[
  \hat u_b =
  \frac{Y_{ab}}{\|Y_{ab}\|_2+\varepsilon},
\]
and the stopped target direction
\[
  u_b =
  \frac{\operatorname{sg}(Z_b)}
       {\|\operatorname{sg}(Z_b)\|_2+\varepsilon}.
\]
The direction loss is
\[
  \mathcal{L}_{\mathrm{dir}}
  =
  1-\mathbb{E}\left[\hat u_b\cdot u_b\right].
\]

The raw target log-magnitude delta is
\[
  \Delta m_{ab}^{\mathrm{raw}}
  =
  \log\left(\|\operatorname{sg}(Z_b)\|_2+\varepsilon\right)
  -
  \log\left(\|\operatorname{sg}(Z_a)\|_2+\varepsilon\right).
\]
With \texttt{--shortcut-mag-scale} equal to $3.0$, the normalized target is
\[
  \Delta m_{ab}^{\star}
  =
  \operatorname{clip}
  \left(
    \frac{\Delta m_{ab}^{\mathrm{raw}}}{3.0},
    -1,1
  \right).
\]
The magnitude loss is
\[
  \mathcal{L}_{\mathrm{mag}}
  =
  \frac{1}{2}
  \mathbb{E}
  \left[
    \left\|
      \Delta m_{ab}
      -
      \operatorname{sg}(\Delta m_{ab}^{\star})
    \right\|_2^2
  \right].
\]

The direct shortcut loss is
\[
  \mathcal{L}_{\mathrm{direct}}
  =
  1.0\,\mathcal{L}_{\mathrm{dir}}
  +
  0.375\,\mathcal{L}_{\mathrm{mag}}.
\]

\section{Bootstrap consistency loss}

The bootstrap branch samples a triplet
\[
  a<b<c
\]
from uniformly sampled positive gaps satisfying $c\le L$.

First, the online predictor maps $a\to b$:
\[
  (Y_{ab},\Delta m_{ab})
  =
  P_\phi(Z_a,a,b,e_\tau,m_a,y).
\]
Then it maps the predicted state $b\to c$:
\[
  (Y_{bc},\Delta m_{bc})
  =
  P_\phi(Y_{ab},b,c,e_\tau,m_b,y).
\]
The EMA predictor gives a stopped one-hop target:
\[
  (Y_{ac}^{\mathrm{ema}},\Delta m_{ac}^{\mathrm{ema}})
  =
  P_{\bar\phi}(\operatorname{sg}(Z_a),a,c,e_\tau,m_a,y).
\]

The direction bootstrap loss is
\[
  \mathcal{L}_{\mathrm{boot,dir}}
  =
  1-
  \mathbb{E}
  \left[
    \frac{Y_{bc}}{\|Y_{bc}\|_2+\varepsilon}
    \cdot
    \operatorname{sg}
    \left(
      \frac{Y_{ac}^{\mathrm{ema}}}
           {\|Y_{ac}^{\mathrm{ema}}\|_2+\varepsilon}
    \right)
  \right].
\]
The magnitude bootstrap loss is
\[
  \mathcal{L}_{\mathrm{boot,mag}}
  =
  \frac{1}{2}
  \mathbb{E}
  \left[
    \left\|
      \Delta m_{ab}+\Delta m_{bc}
      -
      \operatorname{sg}(\Delta m_{ac}^{\mathrm{ema}})
    \right\|_2^2
  \right].
\]
With the command weights,
\[
  \mathcal{L}_{\mathrm{boot}}
  =
  0.25\,\mathcal{L}_{\mathrm{boot,dir}}
  +
  0.1875\,\mathcal{L}_{\mathrm{boot,mag}}.
\]
Because \texttt{--shortcut-bootstrap-detach-source} is set, the bootstrap source
activation and source magnitude are stop-gradient.

\section{Output distillation}

Output distillation is enabled with ratio $0.10$ and weight $0.05$. A subset of
the local batch is sampled, and another centered layer pair $(a_o,b_o)$ is
drawn. The predictor produces
\[
  (Y_{a_ob_o},\Delta m_{a_ob_o})
  =
  P_\phi(Z_{a_o},a_o,b_o,e_\tau,m_{a_o},y).
\]
The predicted target activation is reconstructed as
\[
  \hat Z_{b_o}
  =
  \exp
  \left(
    \operatorname{clip}
    \left(
      m_{a_o}+3.0\,\Delta m_{a_ob_o},
      3.0,8.0
    \right)
  \right)
  \frac{Y_{a_ob_o}}{\|Y_{a_ob_o}\|_2+\varepsilon}.
\]
The DiT tail is resumed from layer $b_o+1$:
\[
  \hat v^{\mathrm{tail}}
  =
  F_\theta^{\mathrm{tail}}
  \left(
    x_\tau,\tau,y;\,
    \mathrm{resume\_hidden}=\hat Z_{b_o},
    \mathrm{resume\_start}=b_o+1
  \right).
\]
The teacher is the stopped full-backbone velocity already computed in the main
forward:
\[
  v^{\mathrm{teacher}}=\operatorname{sg}(\hat v_\theta).
\]
The relative output distillation loss is
\[
  \mathcal{L}_{\mathrm{out}}
  =
  \mathbb{E}
  \left[
    \frac{
      \|\hat v^{\mathrm{tail}}-v^{\mathrm{teacher}}\|_2^2
    }{
      \|v^{\mathrm{teacher}}\|_2^2+10^{-6}
    }
  \right].
\]

The update mode is \texttt{predictor\_plus\_all}. Therefore gradients from this
branch can update the predictor, the DiT tail, and the source-hidden path in the
backbone. The teacher velocity remains stop-gradient.

\section{Private activation loss}

Let post-block activations be
\[
  Z_1,\ldots,Z_L.
\]
With residual private loss enabled, normalize each layer activation and compute
the common component
\[
  C=\frac{1}{L}\sum_{\ell=1}^{L}\operatorname{norm}(Z_\ell).
\]
The private residual is
\[
  R_\ell=\operatorname{norm}(Z_\ell)-\operatorname{sg}(C).
\]
In \texttt{bnd} mode, each $R_\ell$ is flattened over batch, tokens, and
channels. A random set $\mathcal{P}$ of at most four upper-triangular layer pairs
is selected. The loss is
\[
  \mathcal{L}_{\mathrm{private}}
  =
  \frac{1}{|\mathcal{P}|}
  \sum_{(i,j)\in\mathcal{P}}
  \cos^2(R_i,R_j).
\]

\section{Total objective}

The optimized loss for this command is
\[
\boxed{
  \mathcal{L}_{\mathrm{total}}
  =
  \mathcal{L}_{\mathrm{gen}}
  +
  \mathcal{L}_{\mathrm{direct}}
  +
  \mathcal{L}_{\mathrm{boot}}
  +
  0.05\,\mathcal{L}_{\mathrm{out}}
  +
  1.0\,\mathcal{L}_{\mathrm{private}}.
}
\]
Equivalently,
\[
\begin{aligned}
  \mathcal{L}_{\mathrm{total}}
  &=
  \mathcal{L}_{\mathrm{gen}}
  +
  \left(
    1.0\,\mathcal{L}_{\mathrm{dir}}
    +
    0.375\,\mathcal{L}_{\mathrm{mag}}
  \right) \\
  &\quad+
  \left(
    0.25\,\mathcal{L}_{\mathrm{boot,dir}}
    +
    0.1875\,\mathcal{L}_{\mathrm{boot,mag}}
  \right)
  +
  0.05\,\mathcal{L}_{\mathrm{out}}
  +
  \mathcal{L}_{\mathrm{private}}.
\end{aligned}
\]

\section{Training procedure}

Each training step performs:
\begin{enumerate}
\item Load a batch of clean VAE latents and labels $(x_0,y)$.
\item Sample discrete timestep index $q$ from the logit-normal schedule.
\item Convert $q$ to $\tau=q/(T-1)$.
\item Sample noise $x_1\sim\mathcal{N}(0,I)$.
\item Build $x_\tau=(1-\tau)x_1+\tau x_0$.
\item Run a full DiT forward and collect velocity, timestep embedding, and all
  hidden states.
\item Compute $\mathcal{L}_{\mathrm{gen}}$.
\item Sample one centered direct pair and compute
  $\mathcal{L}_{\mathrm{direct}}$.
\item Sample one bootstrap triplet and compute
  $\mathcal{L}_{\mathrm{boot}}$ against the predictor EMA target.
\item Sample an output-distillation subset and centered pair, predict the target
  activation, resume the DiT tail, and compute $\mathcal{L}_{\mathrm{out}}$.
\item Compute the private activation loss on post-block activations.
\item Sum all active losses, backpropagate, and update online backbone and
  predictor parameters.
\item Update the backbone EMA with decay $0.9999$.
\item Update the predictor EMA with the configured predictor EMA decay.
\item Update the L2 magnitude EMA bins for skip-eval magnitude restoration and
  monitoring.
\end{enumerate}

\section{FID sampling}

The command uses
\[
  \texttt{--fid-num-steps}=50,\qquad
  \texttt{--fid-cfg-scale}=1.0,\qquad
  \texttt{--num-fid-samples}=4096.
\]
Sampling starts from Gaussian latent noise:
\[
  x_{\mathrm{init}}\sim\mathcal{N}(0,I).
\]
Class labels are sampled uniformly from ImageNet classes. The sampler runs the
SDE denoising loop in the repo's forward convention from noisy to clean.

\paragraph{Full FID.}
The full sampler uses the EMA DiT at every denoising step:
\[
  \hat v_t = F_{\bar\theta}(x_t,t,y).
\]
The resulting latents are decoded with the VAE, Inception features are computed,
and the run logs
\[
  \texttt{val/FID},\qquad \texttt{val/sFID}.
\]

\paragraph{Depth-shortcut FID.}
The branch also evaluates the depth shortcut because \texttt{--fid-skip-eval}
defaults to on. The shortcut is fixed to
\[
  3\to7.
\]
With
\[
  \texttt{--fid-skip-timestep-mode alternate},
\]
even denoising steps use the shortcut and odd steps use the full DiT.

At a shortcut step, the EMA DiT runs until layer $3$, the predictor EMA maps
\[
  \hat Z_7
  =
  P_{\bar\phi}(Z_3,3,7,e_t,m_3,y),
\]
and the EMA DiT resumes from layer $8$:
\[
  \hat v_t
  =
  F_{\bar\theta}^{\mathrm{tail}}
  \left(
    x_t,t,y;\,
    \mathrm{resume\_hidden}=\hat Z_7,
    \mathrm{resume\_start}=8
  \right).
\]
At a non-shortcut step,
\[
  \hat v_t=F_{\bar\theta}(x_t,t,y).
\]
This reports
\[
  \texttt{val/FID\_skip\_3to7},\qquad
  \texttt{val/sFID\_skip\_3to7}.
\]

\paragraph{Important distinction.}
This command does not enable the timestep-shortcut chain or fanout FID variants.
It evaluates the normal full sampler and the depth shortcut $3\to7$ sampler.

\end{document}
